\subsection{Exact integration of the scale-sensitive triangle parameter}

This reduction starts from the explicitly prepared recursive BPH integrand
of the nested triangle. It does not replace that recursion by a general
forest identity or alter any finite Taylor counterterm. Let
$S=\max_i s_i>0$, $\lambda=S/M$, $L=\log\lambda$, and write all external
invariants and vectors in the following formulas in units of $S$.
For fixed inner and outer simplex fractions put
\[
 \begin{gathered}
 C=1-A,\quad H=u(1-u),\quad D=A+CH,\\
 T=(1-u)z(1-z)\widehat s_1,\quad v=(1-u)z,\quad w=(1-u)(1-z).
 \end{gathered}
\]
The previously established identity gives
\[
 G-T=H\left|k_0-(1-z)\widehat p_1\right|^2=:Z\ge0,
 \qquad
 \delta=\frac{C(AG+CB\widehat s_1)}D,
 \quad B=uvw.
\]
The paired density, including $dA$ but excluding the two simplex measures,
is
\[
 \left\{\frac{\log(1+\lambda AW)}{AC}
 -\frac{A^2}{CD^3}\log(1+\lambda(AW+\delta))\right\}dA.
\]
Its absolute integrability was proved before this change of variables.

Set $x=A/D$, $E=1-(1-H)x$. Then
\[
 A=\frac{Hx}{E},\quad C=\frac{1-x}{E},\quad D=\frac H E,
 \qquad
 \frac{dA}{AC}=\frac{dx}{x(1-x)},\quad
 \frac{A^2dA}{CD^3}=\frac{x^2dx}{1-x}.
\]
Writing $L_c=\log(1+\lambda AW)$ and
$L_r=\log(1+\lambda(AW+\delta))$, the exact $A$ integral is therefore
\[
 I_A=\int_0^1\left\{\frac{L_c}{x}
       +\frac{L_c-L_r}{1-x}+(1+x)L_r\right\}dx.
\]
In particular the cancellation at $x=1$ is retained before integration.

The two positive logarithm arguments are
\begin{align*}
 e^{L_c}&=\frac{1+(K-1)x}{E},\qquad K=H(1+\lambda W),\\
 e^{L_r}&=\frac{N_0(1+rx)(1+sx)}E,\qquad N_0=1+\lambda T,\\
 r+s&=\frac{-(1-H)+\lambda(HW+G-2T)}{N_0},\qquad
 rs=\frac{\lambda(T-G)}{N_0}\le0.
\end{align*}
Thus $r\ge0$, $-1<s\le0$. Evaluation at $x=1$ gives the essential identity
\[
 N_0(1+r)(1+s)=K.
\]
Define the continuous real function
\[
 \mathcal J(t)=\int_0^1(1+x)\log(1+tx)dx
 =\left(\frac32+\frac1t-\frac1{2t^2}\right)\log(1+t)
       -\frac54+\frac1{2t},\quad
 \mathcal J(0)=0,\quad\mathcal J(-1)=-\frac74.
\]
Direct integration of the three finite terms above gives
\begin{align*}
 I_A={}&\frac12\log^2K+\operatorname{Li}_2(1-H)
 +\operatorname{Li}_2\left(\frac r{1+r}\right)
 +\operatorname{Li}_2\left(\frac s{1+s}\right)\\
 &+\frac32\log N_0+\mathcal J(r)+\mathcal J(s)-\mathcal J(H-1).
\end{align*}
Here the first two dilogarithms from the $x^{-1}$ and $(1-x)^{-1}$
integrals combine by
\[
 -\operatorname{Li}_2(1-K)-\operatorname{Li}_2(1-1/K)
 =\frac12\log^2K.
\]
These are exact finite BPH identities, not a fit to high-scale values.

\paragraph{A stable log-domain evaluation.}
Let $R=1+r$, $Q=1+s$, and $\mu=M/S$. Use the bounded scaled quantities
\[
 A_0=\mu+T,\quad B_0=(1+H)\mu+HW+T+Z,\quad K_0=H(\mu+W),
 \quad d_1=-(1-H)\mu+HW+Z-T,
\]
and
\[
 V=\frac{B_0+\operatorname{hypot}(d_1,2\sqrt{A_0Z})}{2},\qquad
 \log R=\log V-\log(\mu+T),\quad \log Q=\log K_0-\log V.
\]
No subtraction of discriminants or construction of a huge root is needed.
The sum $\log(\mu+T)$ uses log-sum-exp, retaining $-L$ when $T=0$.
Landen's identity changes the only negative dilogarithm argument to
\[
 \operatorname{Li}_2\left(\frac{s}{1+s}\right)
 =-\operatorname{Li}_2(1-Q)-\frac12\log^2Q.
\]
All remaining dilogarithm evaluations consequently have arguments in
$[0,1]$. Reflection reduces them to a positive series at most $1/2$; its
tail after $N$ terms is at most
$x^{N+1}/[(N+1)^2(1-x)]$. Small arguments of $\mathcal J$ use its
convergent series rather than cancellation in the displayed closed form.

\paragraph{Separating the large scale before quadrature.}
Write $w_\mu=\log H+\log(\mu+W)$ and
$q_R=\mathcal J(R-1)-3\log R/2$. Algebraically, the exact residual kernel
after removing the explicit scale polynomial is
\begin{align*}
 \mathcal C_\lambda={}&\frac12w_\mu^2
 +L\log(1+\mu/W)+\frac32\log V+\operatorname{Li}_2(1-H)\\
 &+\operatorname{Li}_2(1-R^{-1})-\operatorname{Li}_2(1-Q)
 -\frac12\log^2Q+q_R+\mathcal J(Q-1)-\mathcal J(H-1).
\end{align*}
In particular $q_R\to-5/4$ when $R\to\infty$, so no growing logarithm
remains hidden in that evaluation. Since each simplex has area $1/2$ and
$\int_{\Delta_i}\log H=-1$, the complete amplitude is exactly
\[
 J(\lambda\widehat s)=\frac18L^2+b(\widehat s)L+C_\lambda(\widehat s),
 \qquad
 b=\frac12\int_{\Delta_o}\log W-\frac18,
 \quad C_\lambda=\int_{\Delta_o}\int_{\Delta_i}\mathcal C_\lambda.
\]
The shape moment in $b$ is integrated separately at orders128 and256.
The residual uses scale-independent cubic or quartic Beta-CDF endpoint
maps. Their complements are evaluated directly before binary64 rounding.
The logarithm of $S/M$ is computed from binary mantissas and exponents,
without forming that possibly overflowing ratio.

The independent uniform estimate in the following asymptotic-reference
section bounds $|C_\lambda-c|$ by $130\lambda^{-1/4}$. It controls finite
mass effects when $\mu$ underflows, but is distinct from numerical error in
the shape coefficients or residual quadrature. In the legal zero-outer-leg
shape, $s_3=0$ and $s_1=s_2$, the outer fraction $y$ disappears identically;
the exact-$A$ route then has only three numerical coordinates. Canonical
ordering of the interchangeable outer legs applies this same reduction
when $s_2=0$.
